A level equation $p_U=kp_V$ has centre $\dfrac{(U-kV)}{(1-k)}$. With the first given centre, $a(B-(2,4))=A-(2,4)=(0,-1)$, so B has first coordinate two. The second gives $B=(-1,-2)+b(1,3)$, hence b=3 and B=(2,7). The first equation now gives a=-$\dfrac{1}{3}$. Thus $\dfrac{a}{b}$=-$\dfrac{1}{9}$ and the final centre is $\dfrac{(9A+C)}{10}=\boxed{(\dfrac{9}{5},\dfrac{14}{5})}$.
\begin{asy}
import graph;import geometry;usepackage("amsmath");size(7cm);
pair A=(2,3),B=(2,7),C=(0,1),U=(2,4),V=(-1,-2),W=(1.8,2.8);draw(V--B,dashed);draw(A--B,dashed);dot(A);dot(B);dot(C);dot(U,blue);dot(V,blue);dot(W,red);label("$A$",A,(1,-1));label("$B$",B,(0,1));label("$C$",C,(-1,0));label("$(2,4)$",U,(1,0));label("$(-1,-2)$",V,(-1,0));label("final centre",W,(-1,-1));
\end{asy}

